\documentclass[11pt]{article}
\usepackage[margin=1in]{geometry}
\usepackage{amsmath,amssymb,amsthm,mathtools,microtype,hyperref}
\hypersetup{hidelinks}
\newtheorem{theorem}{Theorem}
\newtheorem{lemma}[theorem]{Lemma}
\newcommand{\ddc}{dd^c}
\newcommand{\Herm}{\operatorname{Herm}}
\newcommand{\tr}{\operatorname{tr}}
\newcommand{\vol}{\operatorname{Vol}}
\newcommand{\Frob}{\mathrm F}
\newcommand{\C}{\mathbb C}
\title{Complete Gauduchon volume geometry for nonclosed parallelizable cones}
\author{}
\date{4 October 2026}
\begin{document}
\maketitle
\begin{abstract}
We describe the intrinsic metric completion of explicit finite-dimensional
cones of nonclosed backgrounds. On a compact quotient of a complex Lie
group, invariant Hermitian forms are balanced and their Gauduchon root
volume is exactly a determinant. We determine the unique normalized
minimizing metric. Every fixed linear subspace of invariant forms containing
a positive form has a complete, proper unit-volume hypersurface for the
negative logarithmic-volume Hessian. Spectral lower bounds and normalized
chord upper bounds give its entire Cauchy classification without enlarging
the prescribed cone. Positive-current coefficient boundary points are
precisely semipositive matrices, but none adds a unit-volume completion
point. We give all constants for the nine-dimensional Iwasawa cone.
The proof uses no bounded-mass hypothesis or Monge--Amp\`ere existence theorem.
\end{abstract}

\section{The specified cones and the completion statement}
Use $d^c=(i/(4\pi))(\bar\partial-\partial)$ and
$\ddc=(i/(2\pi))\partial\bar\partial$.
For a smooth positive probability volume $\mu$ set
\[
 D_\mu(\gamma)=\int_X(\mu/\gamma^n)^{1/n}\gamma^n,\qquad
 \lambda_\mu(\theta)=\inf_{\ddc\gamma^{n-1}=0}
       \frac{\int_X\theta\wedge\gamma^{n-1}}{D_\mu(\gamma)},
\]
where the infimum is over smooth positive Hermitian metrics.
The variational volume is $\lambda_\mu^n$.

Let $G$ be a connected complex Lie group of dimension $n\geq2$ and
$\Gamma$ a discrete subgroup acting on the left with compact quotient
$X=\Gamma\backslash G$. Choose a left-invariant holomorphic coframe
$\alpha_1,\ldots,\alpha_n$, which descends to $X$. For $A\in\Herm(n)$ put
\[
 \theta_A=\frac{i}{2\pi}\sum_{j,k}A_{jk}\alpha_j\wedge\bar\alpha_k,
 \qquad h_0=\theta_I,\qquad H=\int_Xh_0^n,\qquad \mu=h_0^n/H.
\]
Fix a real linear subspace $W\subset\Herm(n)$ containing a positive
matrix and set $V_W=\{\theta_A:A\in W\}$.
The prescribed cone consists of backgrounds in this space admitting a
smooth positive representative after addition of $\ddc u$.
All distances below are intrinsic path distances in that fixed cone or
its unit-volume slice, with metric $g=-D^2\log\lambda_\mu$.

\begin{theorem}\label{thm:main}
For these data the prescribed cone is precisely
$C_W=\{A\in W:A>0\}$ and $V_W\cap\ddc C^\infty(X,\mathbb R)=\{0\}$.
For $A\in C_W$ one has
\begin{equation}\label{eq:volume}
 \lambda_\mu(\theta_A)=(H\det A)^{1/n},\qquad
 g_A(U,Z)=\frac1n\tr(A^{-1}UA^{-1}Z).
\end{equation}
Every positive $\theta_A$ is balanced. The unique minimizing Gauduchon
metric with $D_\mu(\gamma)=1$ is
$\gamma_A=\theta_A/\lambda_\mu(\theta_A)$.

The hypersurface $S_W=\{A\in C_W:H\det A=1\}$ is complete and proper.
Its Cauchy sequences are exactly the sequences converging in coefficients
to a positive matrix in $S_W$. Its completion therefore adds no points,
and every completion point has its canonical smooth positive current
representative $\theta_A$.

A coefficient background $\theta_A$ has a positive-current representative
$\theta_A+\ddc u$, even allowing distributional $u$, exactly when $A\geq0$.
These semipositive coefficient boundary points must not be adjoined to
the unit-volume completion. Nonzero singular ones have infinitely distant
unit-volume normalized minimizing metrics.
If $G$ is nonabelian, every positive background in these cones is nonclosed
and has no closed representative obtained by adding a smooth $\ddc$ potential.
\end{theorem}

The theorem supplies a sufficient-condition classification for a specified
geometric family. It does not classify arbitrary finite-dimensional
nonclosed cones. The probability $\mu$ specified above is part of the
hypotheses.

\section{Compactness, balancedness and the coefficient cone}
Let $Z_j$ be the invariant holomorphic frame dual to $\alpha_j$ and write
$\Omega=\alpha_1\wedge\cdots\wedge\alpha_n$.
Both $\Omega$ and $\bar\Omega$ are closed. For an invariant holomorphic
vector field $Z$, evaluation on the invariant frame gives
\[
 \mathcal L_Z\Omega=-\tr(\operatorname{ad}Z)\Omega,
 \qquad \mathcal L_Z\bar\Omega=0.
\]
The second identity also follows from Cartan's formula, since contraction
of $\bar\Omega$ with $Z$ is zero. The top form $h_0^n$ is a constant
multiple of $\Omega\wedge\bar\Omega$, with the appropriate positive
orientation factor. Stokes on compact $X$ gives
\[
 0=\int_X\mathcal L_Zh_0^n=-H\tr(\operatorname{ad}Z).
\]
Thus the invariant fields have zero divergence and
\begin{equation}\label{eq:mean}
 \int_X Zf\,\mu=\int_X\bar Zf\,\mu=0
 \quad\text{for every smooth complex-valued }f.
\end{equation}
Cartan's formula now yields $d(\iota_Z\Omega)=0$.
These contractions form a basis of the invariant $(n-1,0)$ forms.
Every $\theta_A^{n-1}$ is a constant linear combination of their wedges
with conjugates, so
\begin{equation}\label{eq:balanced}
 d\theta_A^{n-1}=0.
\end{equation}
For $A>0$ this proves balancedness and the Gauduchon condition.

The invariant holomorphic and antiholomorphic fields commute, and
$d\alpha_j$ has type $(2,0)$. The coefficient matrix of $\ddc u$ in the
displayed coframe is therefore $(Z_j\bar Z_k u)$. Its entries have mean
zero by \eqref{eq:mean}. The mean matrix of $\theta_A+\ddc u$ is $A$.
If this form is smooth and positive, its quadratic form on each nonzero
constant coefficient vector is everywhere positive. Integration shows
$A>0$. Conversely $A>0$ supplies the representative $\theta_A$ itself.
If $\theta_A=\ddc u$, taking means gives $A=0$, proving the absence
of nonzero potential directions in $V_W$.

The same mean argument holds for distributions: the formal adjoints
of $Z_j$ and $\bar Z_k$ annihilate constants. The coefficients of a
positive current in this global frame form a positive matrix of measures,
whose mean is semipositive. Thus $\theta_A+\ddc u\geq0$ implies $A\geq0$.
The converse follows from the canonical smooth semipositive form.
When $A=0$, the positive trace measure has total mass zero, and hence
the entire current is zero by polarization. Since $W$ contains a positive
$A_+$, its positive coefficient closure is exactly
$W\cap\{A\geq0\}$: approximate $A\geq0$ by $A+\varepsilon A_+$.

Finally
\[
 \partial\theta_A=\frac{i}{2\pi}\sum_{j,k}A_{jk}
                                  d\alpha_j\wedge\bar\alpha_k.
\]
For invertible $A$, its vanishing forces all $d\alpha_j=0$ by independence
of the conjugate coframe, which means the Lie algebra is abelian.
Thus on a nonabelian quotient every positive invariant background is
nonclosed. Addition of $\ddc u$ cannot alter $d\theta_A$.

\section{Exact volume and the normalized minimizing metric}
The determinant identity gives $\theta_A^n=Q\mu$, where $Q=H\det A$.
For every smooth positive $\gamma$, matrix arithmetic--geometric mean gives
\[
 \frac{\theta_A\wedge\gamma^{n-1}}{\gamma^n}
 =\frac1n\tr_\gamma\theta_A
 \geq(\theta_A^n/\gamma^n)^{1/n}.
\]
After integration,
\[
 \int_X\theta_A\wedge\gamma^{n-1}\geq Q^{1/n}D_\mu(\gamma).
\]
This bounds the quotient for the entire Gauduchon test family.
By \eqref{eq:balanced}, $\gamma=\theta_A$ is an admissible test, with
$D_\mu(\theta_A)=Q^{(n-1)/n}$, and attains equality. This proves
the first formula in \eqref{eq:volume} from the variational definition.

At any minimizing smooth Gauduchon test the continuous nonnegative
arithmetic--geometric mean deficit has integral zero. Pointwise equality
forces $\gamma=f\theta_A$ for a smooth positive function $f$.
Put $q=f^{n-1}$ and $T=\theta_A^{n-1}$. Since $dT=0$, the Gauduchon
equation is $\ddc q\wedge T=0$. Integration by parts gives
\[
 0=\int_Xq\ddc q\wedge T=-\int_Xdq\wedge d^cq\wedge T.
\]
The last integrand is nonnegative and vanishes only where $dq=0$.
Connectedness makes $f$ constant. Since
$D_\mu(c\theta_A)=c^{n-1}\lambda_\mu(\theta_A)^{n-1}$,
normalization $D_\mu=1$ singles out $\gamma_A=\theta_A/\lambda_\mu(\theta_A)$.
On $S_W$ the normalized minimizer is exactly $\theta_A$.

Differentiating the established identity
$\log\lambda_\mu=(\log H+\log\det A)/n$ gives
\[
 D\log\det A[U]=\tr(A^{-1}U),\qquad
 D(A^{-1})[Z]=-A^{-1}ZA^{-1},
\]
and proves the Hessian formula in \eqref{eq:volume}.
For $U\ne0$ its square norm is
$\|A^{-1/2}UA^{-1/2}\|_\Frob^2/n>0$.
It is therefore a smooth positive Riemannian metric on the fixed cone.
Homogeneity makes $S_W$ a smooth hypersurface: the derivative of
$\lambda$ in the radial direction $A$ is $\lambda(A)>0$.

\section{Spectral lower bounds and paths in the fixed cone}
Here $\|\cdot\|_\Frob$ is the Frobenius norm and matrix logarithms are
defined by diagonalization.
\begin{lemma}\label{lem:distance}
On the full positive Hermitian matrix cone with metric \eqref{eq:volume},
\begin{equation}\label{eq:full-distance}
 d_{\rm full}(A,B)=\frac1{\sqrt n}
                   \|\log(A^{-1/2}BA^{-1/2})\|_\Frob.
\end{equation}
The same formula is its intrinsic distance on a full fixed-determinant
hypersurface when $\det A=\det B$.
\end{lemma}
\begin{proof}
Congruence by a fixed invertible matrix is an isometry by the trace
formula. Reduce to $A=I$, $B=C>0$. The path $\exp(t\log C)$ has constant
speed $\|\log C\|_\Frob/\sqrt n$, proving the upper bound.

For the lower bound use the scalar resolvent formula, applied by
diagonalization to $M>0$:
\[
 \log M=\int_0^\infty\bigl((1+s)^{-1}I-(M+sI)^{-1}\bigr)\,ds,
 \quad D\log_M[U]=\int_0^\infty(M+sI)^{-1}U(M+sI)^{-1}\,ds.
\]
The derivative follows by differentiating the inverse; on a neighborhood
with a positive spectral lower bound the derivative integrand has an
integrable bound. In a unitary basis diagonalizing $M$, with eigenvalues
$d_i>0$, this derivative has entries $c_{ij}U_{ij}$, where
\[
 c_{ij}=\frac{\log d_i-\log d_j}{d_i-d_j}\quad(d_i\ne d_j),
 \qquad c_{ij}=1/d_i\quad(d_i=d_j).
\]
Putting $d_i/d_j=e^{2s}$ gives
$\sqrt{d_id_j}|c_{ij}|=|s/\sinh s|\leq1$; the continuous value at
equal eigenvalues gives the same bound. Therefore
\[
 \|D\log_M[U]\|_\Frob^2
 \leq\sum_{i,j}\frac{|U_{ij}|^2}{d_id_j}
 =\|M^{-1/2}UM^{-1/2}\|_\Frob^2.
\]
For any piecewise $C^1$ path from $I$ to $C$, integration and the
Euclidean triangle inequality give
\[
 \operatorname{length}(M)\geq\frac1{\sqrt n}
       \int\left\|\frac{d}{dt}\log M(t)\right\|_\Frob\,dt
 \geq\frac1{\sqrt n}\|\log C\|_\Frob.
\]
This proves \eqref{eq:full-distance}. If $\det A=\det B$, the attaining
path
\begin{equation}\label{eq:geodesic}
 M(t)=A^{1/2}(A^{-1/2}BA^{-1/2})^tA^{1/2}
\end{equation}
has constant determinant, proving the last assertion.
\end{proof}

For a proper $W$, path \eqref{eq:geodesic} need not stay in $W$.
We use the full-matrix formula only for the intrinsic lower bound
\begin{equation}\label{eq:lower}
 d_W(A,B)\geq\frac1{\sqrt n}
                   \|\log(A^{-1/2}BA^{-1/2})\|_\Frob,
 \qquad A,B\in S_W.
\end{equation}
For an upper bound put $U=B-A$, $C(t)=A+tU$ and
$N(t)=C(t)/\lambda(C(t))$.
Convexity of the positive cone and linearity of $W$ put this entire
normalized chord in $S_W$, with endpoints $A,B$.
If $K=C^{-1/2}UC^{-1/2}$, direct differentiation gives
\[
 g_N(N',N')=\frac1n\tr(K^2)-\frac1{n^2}(\tr K)^2
 \leq\frac{\|U\|_\Frob^2}{nm^2},\qquad
 m=\min\{\lambda_{\min}(A),\lambda_{\min}(B)\}>0.
\]
Indeed $C(t)\geq mI$, and the derivative of $\log\lambda(C(t))$
is $\tr(C^{-1}U)/n$. Consequently
\begin{equation}\label{eq:upper}
 d_W(A,B)\leq\frac{\|A-B\|_\Frob}
                  {\sqrt n\min\{\lambda_{\min}(A),\lambda_{\min}(B)\}}.
\end{equation}
The upper bound uses only paths in the prescribed subspace.
In particular $S_W$ is connected.
Writing $d=d_W(A,B)$, \eqref{eq:lower} also yields the matrix order bounds
\begin{equation}\label{eq:order}
 e^{-\sqrt n d}A\leq B\leq e^{\sqrt n d}A.
\end{equation}

\section{All Cauchy sequences and the current boundary}
Let $(A_j)$ be intrinsic Cauchy in $S_W$.
Relative to one fixed tail element, \eqref{eq:order} gives
$mI\leq A_j\leq MI$ for its tail, with fixed $m,M>0$.
Applying the same bounds to two tail elements gives
\[
 \|A_j-A_k\|_{\rm op}\leq M(e^{\sqrt n d_W(A_j,A_k)}-1).
\]
Thus the coefficients are Cauchy. Their limit belongs to the closed
finite-dimensional space $W$, is positive and has $H\det A=1$.
Estimate \eqref{eq:upper} gives intrinsic convergence to that limit.
Conversely coefficient convergence to a positive point of $S_W$ gives
intrinsic convergence by that estimate. This proves the exact Cauchy
criterion and completeness in Theorem~\ref{thm:main}.

The two topologies agree by \eqref{eq:lower}--\eqref{eq:upper}.
A closed distance ball of radius $R$ about $A_*$ lies in the compact
matrix interval
$e^{-\sqrt nR}A_*\leq A\leq e^{\sqrt nR}A_*$, with the closed conditions
$A\in W$ and $H\det A=1$. It is a closed subset of that interval, hence
compact. This proves properness. Distances from $A_*$ are bounded exactly
when the relative eigenvalues are bounded above and away from zero.
Every eigenvalue collapse, explosion or escape from all coefficient
compact subsets of $S_W$ is infinitely distant.

On $S_W$, $\gamma_A=\theta_A$, so \eqref{eq:order} is the pointwise
estimate on the actual normalized minimizing metrics
\[
 e^{-\sqrt n d_W(A,B)}\gamma_A
       \leq\gamma_B\leq e^{\sqrt n d_W(A,B)}\gamma_A.
\]
Finite path length preserves their uniform positivity and boundedness.
Cauchy sequences converge in $C^\infty$, since their coefficients multiply
a fixed smooth coframe. These limits are exactly the canonical positive
current representatives of the already existing points of $S_W$.

For comparison let $A_j\in C_W$ converge in coefficients to a nonzero
singular semipositive matrix $A_0$. Its coefficient boundary has the
canonical current $\theta_{A_0}$, but $\lambda(A_j)\to0$ and
$\widehat A_j=A_j/\lambda(A_j)\in S_W$ has largest eigenvalue tending
to infinity. For fixed $A_*\in S_W$, \eqref{eq:lower} gives, for large $j$,
\begin{align}\label{eq:boundary-rate}
 d_W(A_*,\widehat A_j)
 &\geq\frac1{\sqrt n}\log\frac{\lambda_{\max}(A_j)}
                                  {\lambda(A_j)\lambda_{\max}(A_*)}\\
 &=\frac{|\log\det A_j|}{n\sqrt n}+O(1)\longrightarrow\infty.\notag
\end{align}
The bounded error uses $A_0\ne0$.
Thus positive-current representability at the coefficient boundary
does not imply membership in the unit-volume completion.
At $A_0=0$ normalization may remain interior: $A_j=t_jA_*$ normalizes
to $A_*$. This is radial collapse to an existing shape, not a zero-current
completion point.

The unnormalized cone is infinitely distant from every determinant-zero
boundary. Along every path,
\[
 \left|\frac{d}{dt}\log\lambda(A)\right|
 =\left|\frac1n\tr(A^{-1}A')\right|
 \leq\sqrt{g_A(A',A')}.
\]
This also gives infinite length when volume tends to infinity.
It is the minimizing-metric estimates above that additionally rule out
unit-volume degeneracy.

\section{The full Iwasawa cone with explicit constants}
Take $G=\C^3$ with multiplication
\[
 z\cdot w=(z_1+w_1,z_2+w_2,z_3+w_3+z_1w_2)
\]
and $\Gamma=\mathbb Z[i]^3$. The law and inverse preserve $\Gamma$;
it is discrete. Its left action is free and properly discontinuous:
for compact $K\subset G$, elements with $\gamma K\cap K\ne\varnothing$
belong to the compact set $KK^{-1}$ and there are only finitely many.
Each coset has a representative in the closed real six-dimensional unit
cube. Choose lattice coordinates $\gamma_1,\gamma_2$ to reduce $z_1,z_2$,
then $\gamma_3$ to reduce $z_3+\gamma_1z_2$.
Representatives in the half-open cube are unique, first in their first
two coordinates and then in the third. Thus the cube is a measurable
fundamental domain; its compact image proves compactness of $X$.

The global invariant holomorphic coframe is
\[
 \alpha_1=dz_1,\quad\alpha_2=dz_2,\quad
 \alpha_3=dz_3-z_1dz_2,\qquad
 d\alpha_3=-\alpha_1\wedge\alpha_2.
\]
Set $W=\Herm(3)$, of real dimension nine.
Since $\alpha_1\wedge\alpha_2\wedge\alpha_3=dz_1\wedge dz_2\wedge dz_3$,
the fundamental-domain density of $h_0^3$ is $6/\pi^3$ times real
Lebesgue volume. Therefore
\[
 H=6/\pi^3,\qquad \lambda_\mu(\theta_A)^3=(6/\pi^3)\det A,\qquad
 S=\{A>0:\det A=\pi^3/6\}.
\]
Every positive $\theta_A$ is genuinely nonclosed: the coefficient of
$\alpha_1\wedge\alpha_2\wedge\bar\alpha_3$ in $\partial\theta_A$
is $-iA_{33}/(2\pi)$ and $A_{33}>0$.

The quotient is also non-K\"ahler, which can be checked directly.
If a positive closed metric $\kappa$ existed, write its coefficient
functions in the same coframe and average them against $\mu$.
This gives a positive matrix $A$. In the averaged coefficients of
$\partial\kappa$, all derivative terms integrate to zero by
\eqref{eq:mean}, so their mean is $\partial\theta_A$.
Closedness would force $\partial\theta_A=0$, contradicting $A_{33}>0$.
No classification theorem is needed for this assertion.

Here Lemma~\ref{lem:distance} gives the exact intrinsic distance on the
eight-dimensional unit-volume slice:
\[
 d_S(A,B)=\frac1{\sqrt3}\|\log(A^{-1/2}BA^{-1/2})\|_\Frob.
\]
Path \eqref{eq:geodesic} stays in this prescribed full cone and in $S$.
For $A_*=aI$, $a=(\pi^3/6)^{1/3}$,
\[
 d_S(A_*,A)^2=\frac13\sum_{i=1}^3
                          \left(\log\frac{\lambda_i(A)}a\right)^2.
\]
The three logarithms sum to zero. Every genuine spectral degeneration
has infinite distance; every Cauchy sequence converges to a smooth
positive invariant minimizer. Nonzero semipositive coefficient boundary
currents have normalized distance at least
$|\log\det A_j|/(3\sqrt3)+O(1)$ and add no completion point.

The same fixed-cone classification applies to every separate real linear
subspace $W$ on every compact quotient in Theorem~\ref{thm:main}.
For proper $W$ we retain the spectral lower bound and the intrinsic
normalized-chord upper bound, rather than asserting a geodesic formula
that could leave $W$. The conclusion concerns these invariant-volume
cones only; other geometries can have finite-distance nef boundaries.
\end{document}
